%!TEX root = ../main.tex

% ==================================================================
\section{Lecture~1: The why and how of Calabi--Yau}\label{sec:L1}

\medskip

\noindent\textit{Anchor:} TASI \cite{McAllister:2025qwq},
\S\S\,1--2.

\medskip
\noindent\textit{Guiding question.} What can toric methods compute
for string compactifications, and where do current methods still
fall short?

The answer will be deliberately two-sided. Toric geometry gives an
explicit, computable supply of Calabi--Yau threefold hypersurfaces:
Hodge numbers, intersection forms, cones, divisor topology, periods,
and enumerative inputs such as Gopakumar--Vafa invariants can be
extracted algorithmically. It does not, by itself, solve the
Ricci-flat metric, determine every effective curve or divisor on the
hypersurface, or provide access to all quantum corrections.

\medskip
\noindent\textit{Historical and computational thread.}
The geometric problem starts with two classical inputs. Yau's
solution of the Calabi conjecture implies the existence of
Ricci-flat K\"ahler metrics on compact K\"ahler manifolds with
trivial canonical bundle \cite{Yau:1978cfy}, while Calabi--Yau
threefolds entered string compactification through the work of
Candelas, Horowitz, Strominger and Witten
\cite{Candelas:1985en}. The reason these lectures focus on toric
hypersurfaces is more computational: Batyrev's reflexive polytopes
give a large, explicit class of mirror pairs
\cite{Batyrev:1993oya}; the Kreuzer--Skarke classification makes
this class finite and searchable \cite{Kreuzer:2000xy}; and modern
packages such as \texttt{CYTools} turn the resulting combinatorial
data into the intersection, cone and divisor-topology inputs used
later \cite{Demirtas:2022hqf}. Lecture~1 explains enough of this
chain to make Lectures~2 and~3 self-contained.

\medskip
\noindent\textit{Goals.} After this lecture you can
\begin{enumerate}[nosep,leftmargin=*]
    \item state what a string compactification is, as a geometric
    object, and why the vacuum Einstein equations on a product ansatz
    select Ricci-flat internal manifolds;
    \item explain why Calabi--Yau threefolds are the computable
    Ricci-flat class used in these lectures, and how reduced holonomy
    gives the parallel spinors responsible for residual
    supersymmetry;
    \item identify the moduli of a Calabi--Yau compactification with
    specific cohomological data;
    \item explain why Calabi--Yau hypersurfaces in toric varieties
    are the natural constructive source of explicit compact examples,
    and what data they feed into Lectures~2 and~3.
\end{enumerate}


\subsection{Why string theory, and why compactify?}



Let us motivate how geometry enters physics.
We have very successful but structurally different theories
for matter and gravity:
\begin{equation}
\begin{array}{lll}
\text{matter} &:&
  \begin{array}[t]{l}
    \text{Standard Model},\\
    \text{\emph{language:} quantum field theory},\\
    \text{\emph{hypothesis:} point particles};
  \end{array}
\\[1.2em]
\text{gravity} &:&
  \begin{array}[t]{l}
    \text{General Relativity},\\
    \text{\emph{language:} differential geometry},\\
    \text{\emph{hypothesis:} spacetime as a manifold}.
  \end{array}
\end{array}
\end{equation}
The quantum gravity problem is to find a framework in which these two
descriptions can be combined.  Perturbative quantum gravity based on
point particles is not perturbatively renormalisable, and curvature
singularities indicate that the classical spacetime description
itself can break down.  String theory is one proposed framework in
which the fundamental extended objects are one-dimensional strings
rather than point particles.

\begin{figure}
\begin{center}
\centering
\begin{tikzpicture}[scale=1.08,>=Latex]
  \begin{scope}
    \draw[->,thin] (-0.25,-0.15) -- (2.35,-0.15) node[right] {$x$};
    \draw[->,thin] (-0.25,-0.15) -- (-0.25,2.70) node[above] {$t$};
    \draw[thick,->] (0.55,0.15)
      .. controls (0.35,0.70) and (0.95,1.05) .. (0.75,1.55)
      .. controls (0.60,1.95) and (1.10,2.30) .. (1.35,2.52);
    \fill (0.55,0.15) circle (2pt);
    \fill (0.75,1.55) circle (2pt);
    \node[align=center] at (1.05,-0.78) {particle theory:\\ worldline};
  \end{scope}

  \node at (3.35,1.25) {\large vs.};

  \begin{scope}[xshift=4.55cm]
    \draw[->,thin] (-0.25,-0.15) -- (2.75,-0.15) node[right] {$x$};
    \draw[->,thin] (-0.25,-0.15) -- (-0.25,2.70) node[above] {$t$};
    \path[fill=blue!6]
      (0.55,0.20)
      .. controls (0.30,0.85) and (0.78,1.45) .. (0.55,2.55)
      -- (1.65,2.55)
      .. controls (1.95,1.78) and (1.30,1.00) .. (1.58,0.20)
      -- cycle;
    \begin{scope}
      \clip
        (0.55,0.20)
        .. controls (0.30,0.85) and (0.78,1.45) .. (0.55,2.55)
        -- (1.65,2.55)
        .. controls (1.95,1.78) and (1.30,1.00) .. (1.58,0.20)
        -- cycle;
      \draw[thin] (0.15,0.72)
        .. controls (0.65,0.55) and (1.18,0.90) .. (2.05,0.74);
      \draw[thin] (0.15,1.42)
        .. controls (0.70,1.25) and (1.24,1.62) .. (2.05,1.45);
      \draw[thin] (0.15,2.12)
        .. controls (0.70,1.95) and (1.28,2.30) .. (2.05,2.15);
    \end{scope}
    \path[draw,thick]
      (0.55,0.20)
      .. controls (0.30,0.85) and (0.78,1.45) .. (0.55,2.55)
      -- (1.65,2.55)
      .. controls (1.95,1.78) and (1.30,1.00) .. (1.58,0.20)
      -- cycle;
    \draw[->,thin] (2.25,0.35) -- (2.25,2.35);
    \node[align=center] at (1.10,-0.78) {string theory:\\ worldsheet $\Sigma$};
  \end{scope}
\end{tikzpicture}
\end{center}
\vspace*{-0.5cm}
\caption{A point particle sweeps out a
one-dimensional worldline, whereas a string sweeps out a
two-dimensional worldsheet.  In both panels the object evolves upward
in target time $t$: the point traces a curve, while the spatially
extended string traces a strip. }
\end{figure}

The useful contrast is the following. A particle
traces out a worldline in spacetime; a string traces out a worldsheet
$\Sigma$, and the simplest classical action is proportional to its
area,
\begin{equation}
    S_{\rm NG}=-T\int_{\Sigma}\mathrm{d}A \, ,
\end{equation}
where $T$ is the string tension. Quantising the string gives a
tower of vibrational modes. 
In the point-particle picture
one quantises the possible motions of a point; in the string picture
one also quantises the normal modes of an extended object.  Very
schematically,
\begin{equation}
    \text{particle: worldline}
    \qquad\longrightarrow\qquad
    \text{string: worldsheet with oscillator modes}\, .
\end{equation}
The word ``harmonics'' here refers to these oscillator modes of the
string.  After quantisation they appear, from the spacetime point of
view and for an observer sufficiently far away, as different particle
species: scalars, gauge fields, spinors, and a spin-$2$ graviton
among them.


It is useful to separate two attitudes.  One can use
string theory and related constructions as mathematical
tools in broader effective descriptions.  In these lectures we use
the stronger quantum gravity interpretation: the theory is required
to include gravity, and the perturbative spectrum contains a
spin-$2$ graviton.  With this interpretation, the consistency
conditions of the quantum theory are part of the input.  In
particular, for the critical superstring the consistency of the quantum theory fixes the target spacetime dimension to ten.
At energies small compared with the
string scale, one keeps only the massless modes. 
The massive oscillator modes have masses of order the string
scale $M_s$, so a low-energy observer with characteristic energy
$E\ll M_s$ cannot excite them as dynamical fields.  Their effects are
still present indirectly, through higher-derivative corrections to
the effective action.
The one fact from
this construction that Lecture~1 needs is that the massless spectrum
of a consistent superstring includes a ten-dimensional metric,
together with further geometric fields such as differential forms,
scalars, spinors, and gauge fields. The metric is what will make
Ricci-flat geometry enter immediately.

Since observed physics is effectively
four-dimensional at accessible energies, a compactification must
explain how a ten-dimensional theory can give a four-dimensional
effective theory:
\begin{equation}
    \dim \mathcal{M}_{10}=10=9+1,
    \qquad
    \dim \mathcal{M}_4=4=3+1,
    \qquad
    \mathcal{M}_{10}\simeq \mathcal{M}_4\times X .
\end{equation}
The compact six-dimensional factor $X$ is where geometry enters.
The central point of the lectures is that the topology and geometry
of $X$ determine the light four-dimensional fields, their
couplings, and the scalar potential whose critical values are the
vacuum energies.




\subsection{Framing: what a string compactification is}\label{ssec:L1-framing}

For the first lecture, the input from string theory can be stated
without choosing a specific string theory. It is the following
geometric package.

\begin{enumerate}[leftmargin=*]
    \item A \emph{ten-dimensional spacetime} $\mathcal{M}_{10}$:
    a Lorentzian $10$-manifold with metric $g_{MN}$.
    \item Additional \emph{fields} $\Phi$ on $\mathcal{M}_{10}$.
    Here a field simply means a geometric object varying over
    spacetime: a section of a vector bundle, such as a differential
    form, a function, a connection, a spinor, or the metric itself.
    In Lecture~1 we keep these fields mostly silent and set them to zero or
    to constants when deriving the Ricci-flatness condition.
    \item \emph{Source terms and charged objects}, such as branes
    and orientifold planes,
    which enter through $\mathcal{L}_{\rm source}$ and produce
    localised stress-energy and charge.
    A brane is an extended object on which open strings can
    end, or more generally an object carrying charge under one of the
    differential-form gauge fields of the theory.  Its worldvolume is
    a submanifold of spacetime, so after compactification a brane can
    wrap a cycle in $X$ and fill some or all of the visible
    four-dimensional spacetime.  This is why branes are simultaneously
    geometric objects and physical sources: they contribute localised
    energy-momentum, carry quantised charge, and can support gauge
    fields or matter in the four-dimensional theory
    \cite{Polchinski:1996na}.
    \item An \emph{action functional}
    \begin{equation}
        S[g,\Phi]
        =
        \frac{1}{2\kappa_{10}^2}\int_{\mathcal{M}_{10}}
        \mathrm{d}^{10}x\,\sqrt{-g}\,
        \bigl(R^{(10)}+\mathcal{L}_{\rm matter}(g,\Phi)
        +\mathcal{L}_{\rm source}(g,\Phi)+\underbrace{\ldots}_{\mathrm{e.g.}\, R^4}\bigr)\, ,
    \end{equation}
    whose Euler--Lagrange equations are the equations of motion,
    i.e., the differential equations obtained by requiring the first
    variation of $S$ to vanish, see Eq.~\eqref{eq:EOMRvac} below.
    This is the classical, leading-order ten-dimensional action.
    The full string effective action also contains
    higher-derivative and quantum corrections.\footnote{The first well-known perturbative
    correction to the ten-dimensional type-II action contains an
    $R^4$ term at order $(\alpha')^3$. The
    parameter $\alpha'$ is the inverse string tension, up to the
    conventional factor $2\pi$, and sets the square of the string
    length.  Corrections in powers of $\alpha'$ are higher-derivative
    or finite-string-size corrections.  The string coupling $g_s$
    controls the loop expansion.}
    We do not include those corrections in the first-pass
    vacuum equation, but they reappear in Lecture~2 as corrections to
    the four-dimensional effective action.
\end{enumerate}

Let $x^\mu$ be coordinates on a four-dimensional Lorentzian manifold
$\mathcal{M}_4$ and $y^m$ coordinates on a compact real
$6$-manifold $X$. A warped product metric on
$\mathcal{M}_{10}=\mathcal{M}_4\times X$ has the form
\begin{equation}\label{eq:L1-warped-ansatz}
    \mathrm{d}s^2_{10}
    =
    \mathrm{e}^{2A(y)}\,g^{(4)}_{\mu\nu}(x)\,\mathrm{d}x^\mu \mathrm{d}x^\nu
    +
    \mathrm{e}^{-2A(y)}\,g^{(6)}_{mn}(y)\,\mathrm{d}y^m \mathrm{d}y^n .
\end{equation}
The function $A(y)$ is the \emph{warp factor}. In the first pass
through the vacuum equations we set $A=0$, so the metric is the
ordinary product metric
$\mathrm{d}s^2_{10}=\mathrm{d}s^2(\mathcal{M}_4)+\mathrm{d}s^2(X)$.
This is the Kaluza--Klein bridge in its simplest form: a
four-dimensional observer does not see the internal coordinates
directly, but sees the coefficients of expansions in functions,
forms, and metric deformations on $X$.

\begin{definition}[String compactification, geometrically]
A \emph{string compactification} consists of:
(i) a factorisation $\mathcal{M}_{10} = \mathcal{M}_4\times X$
of the ten-dimensional spacetime into a four-dimensional Lorentzian
manifold $\mathcal{M}_4$ and a compact $6$-dimensional Riemannian
manifold $X$ (the \emph{internal} space);
(ii) a solution of the ten-dimensional equations of motion on
$\mathcal{M}_{10}$ with metric in the warped product form
\eqref{eq:L1-warped-ansatz};
(iii) background choices for any non-metric fields and charged
objects that are present, subject to the global consistency
conditions of the chosen string theory.
\end{definition}

On a compact internal space, closed differential-form field
strengths can have quantised periods. Abstractly, if $F_p$ is a
closed $p$-form field strength, \emph{flux quantisation} means that its
normalised periods over integral $p$-cycles are integers, so the
background defines an integral cohomology class. The precise
normalisation and even the correct cohomology theory depend on the
string theory and on the charged objects present. Lecture~2
specialises this statement to Type IIB, where the flux
superpotential uses the three-form classes
$[F_3],[H_3]\in H^3(X,\mathbb{Z})$.

At this point the lecture has only defined the geometric problem: choose a compact six-manifold $X$, a metric on $\mathcal{M}_4\times X$, and background fields and sources
so that the ten-dimensional equations of motion are satisfied, and
then determine the four-dimensional fields and interactions obtained
by expanding around that solution.
We now first find the internal geometries that can solve the equations of motion in the compactification background, and then ask which four-dimensional fields and
potentials those geometries produce.

\subsection{The vacuum solution problem}\label{ssec:L1-vacuum}

We begin in the simplest setting: no fluxes, no branes. We ask which
Riemannian metrics on $X$ are compatible with a vacuum solution.

The metric equation comes from varying the schematic action above
with respect to $g^{MN}$. Let
\begin{equation}
    S_{\rm matter+source}
    =
    \frac{1}{2\kappa_{10}^2}
    \int \mathrm{d}^{10}x\,\sqrt{-g}\,
    \bigl(\mathcal{L}_{\rm matter}+\mathcal{L}_{\rm source}\bigr)\, ,
\end{equation}
where $\mathcal{L}_{\rm matter}$ contains the kinetic terms of the
non-metric fields $\Phi$ and $\mathcal{L}_{\rm source}$ contains
localised brane/orientifold terms. Define
\begin{equation}
    T_{MN}
    \coloneqq  
    -\frac{2}{\sqrt{-g}}\,
    \frac{\delta S_{\rm matter+source}}{\delta g^{MN}} .
\end{equation}
Then the ten-dimensional Einstein equation is
\begin{equation}\label{eq:EOMRvac} 
    R^{(10)}_{MN}
    =
    \kappa_{10}^2
    \left(T_{MN}-\frac18 g_{MN}T^K{}_K\right).
\end{equation}
In this first pass we set all additional fields to zero or constant
and include no localised sources, so $T_{MN}=0$. The unwarped product
version of \eqref{eq:L1-warped-ansatz} then gives
\begin{equation}
    R_{\mu\nu}^{(4)} = 0 \,, \qquad R_{mn}^{(6)} = 0 \,.
\end{equation}
The four-dimensional Einstein vacuum equation forces $\mathcal{M}_4$
to be Ricci-flat. For our purposes the physically interesting case
is $\mathcal{M}_4 = \mathbb{R}^{1,3}$ (Minkowski), so
$R_{\mu\nu}^{(4)}=0$ holds trivially. The interesting condition is
the internal one: \emph{$X$ must be Ricci-flat}.

The problem is then: construct compact Ricci-flat Riemannian
$6$-manifolds on which one can compute. This is a hard nonlinear PDE
problem: explicit Ricci-flat metrics on compact non-flat
$6$-manifolds are generally not known in closed form. The useful
insight is that one does not need the metric explicitly in order to
build the compactification. Yau's theorem guarantees the
\emph{existence} of Ricci-flat K\"ahler metrics from
complex-geometric data, and toric geometry gives a large explicit
class of spaces for which the relevant topological data are
computable.

\subsection{Berger's holonomy classification}\label{ssec:L1-berger}

The compactification problem above asks for Ricci-flat internal
manifolds with enough structure to control the low-energy theory.
Holonomy is the efficient invariant for this purpose. It records how
tangent vectors, tensors and spinors transform under parallel
transport, and therefore tells us which parallel geometric objects a
metric can admit. 

\begin{theorem}[Berger, holonomy classification \cite{Joyce:2000,Besse:1987}]\label{thm:berger}
Let $(X,g)$ be a simply-connected, complete Riemannian manifold
whose holonomy representation is irreducible and whose metric is
not locally symmetric. Then the (restricted) holonomy group
$\Hol(X,g)\subseteq\mathrm{SO}(\dim X)$ is one
of groups in Table~\ref{tab:berger}.
The lower bounds $m\geq 2$ remove degenerate or duplicated
low-dimensional cases ($\mathrm{U}(1)=\mathrm{SO}(2)$,
$\mathrm{Sp}(1)=\mathrm{SU}(2)$,
$\mathrm{Sp}(1)\!\cdot\!\mathrm{Sp}(1)=\mathrm{SO}(4)$).
\end{theorem}

\begin{table}[t!]
\begin{center}
\begin{tabular}{llll}
\toprule
Holonomy & Dimension & Ricci-flat? & Geometry relevant here \\
\midrule
$\mathrm{SO}(n)$ & $n$ & no generic implication & no parallel spinor in general \\
$\mathrm{U}(m)$ & $2m$ & not necessarily & K\"ahler \\
$\mathrm{SU}(m)$ & $2m$ & yes & Calabi--Yau (Kähler+Ricci-flat) \\
$\mathrm{Sp}(m)$ & $4m$ & yes & hyperk\"ahler \\
$\mathrm{Sp}(m)\!\cdot\!\mathrm{Sp}(1)$ & $4m$ & no (Einstein) & quaternionic K\"ahler \\
$G_2$ & $7$ & yes & $G_2$ manifolds \\
$\mathrm{Spin}(7)$ & $8$ & yes & $\mathrm{Spin}(7)$ manifolds \\
\bottomrule
\end{tabular}
\end{center}
\caption{Berger's list of irreducible
Riemannian holonomies.}\label{tab:berger} 
\end{table}


Berger's theorem is useful here because it sharply limits the
possible irreducible, non-symmetric Riemannian holonomies. It does
\emph{not} classify all compact Ricci-flat manifolds. Rather, it
explains why, once one asks for a six-dimensional irreducible
Ricci-flat background admitting parallel spinors, the natural
holonomy group is contained in $\mathrm{SU}(3)$. In the terminology
of \cite{McAllister:2025qwq}, the compact Ricci-flat
examples used in string compactification all arise from reduced, or
special, holonomy; the Calabi--Yau case is the six-dimensional
member of this list that is complex, K\"ahler, and computationally
accessible. This is a guide to the known constructive landscape, not
a theorem excluding every possible Ricci-flat six-manifold outside
the classes we discuss. 



For a compact K\"ahler threefold, the algebraic-geometric condition
that implements the $\mathrm{SU}(3)$ holonomy target is the
triviality of the canonical bundle.
Let $\mathcal{O}_X$ denote the structure sheaf of holomorphic
functions on $X$; an isomorphism $K_X\cong\mathcal{O}_X$ means that
the canonical bundle is holomorphically trivial, equivalently that
there is a nowhere-vanishing holomorphic $(3,0)$-form.
Holomorphic triviality of $K_X$ gives $c_1(X)=0$;
Yau's theorem then gives a Ricci-flat K\"ahler metric in each
K\"ahler class; and, after excluding flat or product factors, the
holonomy is the expected subgroup of $\mathrm{SU}(3)$,
\begin{equation}
    K_X \cong \mathcal{O}_X
    \quad\Longrightarrow\quad
    \text{Ricci-flat K\"ahler metric}
    \quad\Longrightarrow\quad
    \Hol(X)\subseteq \mathrm{SU}(3)\, .
\end{equation}
When the holonomy is exactly $\mathrm{SU}(3)$, this is the generic
irreducible Calabi--Yau threefold case. Proper subgroups correspond
to extra structure, such as flat factors or product geometries, and
usually to more supersymmetry than the generic threefold case. This
is the precise sense in which Calabi--Yau threefolds are the
workhorse Ricci-flat backgrounds of the lectures: they are explicit
enough to construct and compute with, and their reduced holonomy has
the spinorial consequence needed by the physics.

\subsection{Aside: Reduced holonomy, parallel spinors, and supersymmetry\texorpdfstring{$^{\ast}$}{*}}
\label{ssec:L1-ricci-flat}

Supersymmetry is not the starting point of the geometric argument; it
is a consequence of the reduced holonomy just described. The physics
reason one cares about it is that residual supersymmetry gives
technical control of the low-energy theory and organises the
four-dimensional fields into multiplets. The mathematical statement
is the existence of parallel sections of the spinor bundle.

\begin{definition}[Parallel (covariantly constant) spinor]
Let $(X,g)$ be a Riemannian spin manifold with spinor bundle $S\to X$
and let $\nabla$ denote the spin connection. A \emph{parallel} (or
\emph{covariantly constant}) \emph{spinor} is a smooth section
$\eta\in\Gamma(S)$ with $\nabla \eta = 0$.
\end{definition}

\begin{remark}[Terminology]
In the physics literature these are often called \emph{Killing
spinors}. In the standard differential-geometry convention,
a \emph{Killing spinor} satisfies
$\nabla_X\eta = \lambda\, X\!\cdot\!\eta$ for a nonzero constant
$\lambda$, and the $\lambda=0$ case is called a \emph{parallel}
spinor and treated separately. For the unwarped Minkowski
compactification considered here the relevant case is $\lambda=0$,
and we use the term \emph{parallel} throughout.
\end{remark}

\begin{proposition}[Physics dictionary; standard supergravity
    Killing-spinor analysis,
    \cite{McAllister:2025qwq} \S\,2.3]\label{thm:susy-parallel}
For the configuration $\mathcal{M}_4 \times X$ with
$\mathcal{M}_4 = \mathbb{R}^{1,3}$, the unwarped product metric,
no fluxes, and no branes, vanishing of the supergravity SUSY
variations on the four-dimensional vacuum requires that $X$ admit
a parallel spinor. Starting from a ten-dimensional Type II theory
with thirty-two real supercharges, one independent complex parallel
spinor on $X$ corresponds to $\mathcal{N}=2$ four-dimensional
supersymmetry (eight real supercharges) before any orientifold
projection; starting instead from a heterotic or Type I theory
(sixteen real supercharges in ten dimensions), the same condition
yields $\mathcal{N}=1$ in four dimensions (four real supercharges),
as in Candelas--Horowitz--Strominger--Witten
\cite{Candelas:1985en}.

\emph{Status:} this is a standard physics derivation, not a
mathematical theorem about $(X,g)$ alone, and assumes the
leading-order supergravity expansion.
\end{proposition}

The invariant internal spinors determine the residual supersymmetry.
The precise conversion between internal spinor multiplicities and
four-dimensional supercharges depends on the ten-dimensional
chirality and real/complex spinor conventions. We only use the
standard Type IIB case: an SU($3$)-holonomy Calabi--Yau threefold
gives $\mathcal{N}=2$ supersymmetry in four dimensions (eight real
supercharges) before the orientifold projection in Lecture~2 reduces
it to $\mathcal{N}=1$.

The bridge to topology is Berger's classification of holonomies on
Riemannian manifolds. The curvature condition comes from the
commutator of spin covariant derivatives. In an orthonormal frame,
\begin{equation}
    \nabla_m\eta
    = \partial_m\eta
        + \frac14\,\omega_m{}^{ab}\Gamma_{ab}\eta ,
    \qquad
    [\nabla_m,\nabla_n]\eta
    = \frac14\,R_{mn}{}^{ab}\Gamma_{ab}\eta .
\end{equation}
If $\eta$ is parallel, the left-hand side vanishes, hence
\begin{equation}
    R_{mn}{}^{ab}\Gamma_{ab}\eta = 0 .
\end{equation}
By the Ambrose--Singer theorem the restricted holonomy algebra is
generated by curvature endomorphisms, so this integrability
condition says precisely that the holonomy representation fixes a
non-zero spinor. This can happen on a non-flat manifold only if the
holonomy group acts trivially on some subspace of the spinor
representation; see \cite{LawsonMichelsohn:1989} for the spin
geometry and \cite{McAllister:2025qwq}, \S\,2.3 for the physics
dictionary.


\subsection{Calabi--Yau threefolds: existence and Hodge structure}\label{ssec:L1-CY}

The six-dimensional geometry problem has now become concrete: find
compact Riemannian $6$-manifolds whose holonomy is contained in
$\mathrm{SU}(3)$, and preferably equals $\mathrm{SU}(3)$ after
excluding torus and product factors. The systematic way to produce
such spaces is to construct compact K\"ahler threefolds with
trivial canonical bundle and then apply Yau's theorem to obtain
Ricci-flat metrics. Toric Calabi--Yau hypersurfaces are not meant to
exhaust the known compact examples with
$\mathrm{SU}(3)$ holonomy. Rather, Calabi--Yau threefolds are the
standard geometric class relevant to this compactification problem,
and Batyrev hypersurfaces are the large computable family used in
these lectures.

\begin{definition}[Calabi--Yau threefold]\label{def:CY3}
A \emph{Calabi--Yau threefold} is a compact K\"ahler manifold $X$ of
complex dimension $3$ with trivial canonical bundle
$K_X\cong\mathcal{O}_X$ and $\pi_1(X)=0$.
\end{definition}

Equivalently, $X$ admits a nowhere-vanishing holomorphic
$(3,0)$-form $\Omega$, unique up to scale. The appendix records the
standard comparison with other conventions, including formulations in
terms of $c_1(X)$ and holonomy.

The simple-connectedness clause excludes $T^6$, $K3\times T^2$,
abelian threefolds, and CY3s obtained as free quotients by a
finite group with nontrivial $\pi_1$ (e.g.\ the heterotic
Standard-Model quotients of Donagi--Ovrut \cite{DonagiOvrut:2000}
and Bouchard--Donagi \cite{BouchardDonagi:2005}). It implies
$h^{1,0}(X)=0$, removing spurious vector multiplets from
$H^{1,0}$ and giving a clean four-dimensional $\mathcal{N}=1$
truncation in Lecture~2. The converse fails in general: free
finite quotients can have torsion $\pi_1$ while still satisfying
$h^{1,0}=0$. The rigorous classification of compact K\"ahler
manifolds with $c_1(X)=0$ in $H^2(X,\mathbb{R})$ is given by the
Beauville--Bogomolov decomposition \cite{Beauville:1983,Joyce:2000}:
such an $X$ admits a finite \'etale cover that is a product of
complex tori, strict Calabi--Yau factors (with
$\mathrm{Hol}=\mathrm{SU}(n)$, $n\geq 3$), and irreducible
hyperk\"ahler factors (with $\mathrm{Hol}=\mathrm{Sp}(n)$).

\begin{theorem}[Yau \cite{Yau:1978cfy}]\label{thm:yau}
Let $X$ be a compact K\"ahler manifold with $c_1(X)=0$. Then in each
K\"ahler class $[J]\in H^{1,1}(X)\cap H^2(X,\mathbb{R})$ there exists
a unique K\"ahler metric $g$ with $\mathrm{Ric}(g)=0$.
\end{theorem}

\begin{corollary}\label{cor:CY-Ricci}
Calabi--Yau threefolds admit Ricci-flat K\"ahler metrics; in
particular they solve the vacuum Einstein equation
$R_{mn}^{(6)} = 0$ required by the unwarped product ansatz.
\end{corollary}

What follows is the cohomological data we will need in Lecture~2:

\begin{proposition}[CY3 Hodge diamond]
Let $X$ be a Calabi--Yau threefold. The only independent Hodge numbers are $h^{1,1}$ and $h^{2,1}$, and the Euler characteristic is $\chi(X) = 2\bigl(h^{1,1}-h^{2,1}\bigr)$.
\end{proposition}

The holomorphic $(3,0)$-form $\Omega$ is unique up to overall
rescaling and is nowhere vanishing. The pair
$(h^{1,1},h^{2,1})$ is the principal topological invariant of $X$ for
our purposes; we will see that $h^{1,1}$ counts \emph{K\"ahler
moduli} (deformations of $J$) and $h^{2,1}$ counts \emph{complex
structure moduli} (infinitesimal deformations of the complex
structure on $X$, equivalently classes in $H^1(X, T_X)$).

\begin{center}
\begin{tabular}{lll}
\toprule
Geometric datum & Cohomology group & Dimension \\
\midrule
K\"ahler deformations & $H^{1,1}(X,\mathbb{R})$ & $h^{1,1}$ \\
Complex structure deformations & $H^1(X,T_X)\cong H^{2,1}(X)$ & $h^{2,1}$ \\
Holomorphic volume form & $H^{3,0}(X)$ & $1$ \\
\bottomrule
\end{tabular}
\end{center}

\subsection{Massless spectrum and moduli of CY compactifications}\label{ssec:L1-moduli}

We now identify the four-dimensional moduli with cohomological data
on $X$, following \cite{McAllister:2025qwq} \S\,2.5.

\subsubsection{Geometric moduli: K\"ahler and complex structure}

\begin{proposition}[K\"ahler moduli]\label{prop:K-moduli}
The space of infinitesimal Ricci-flat K\"ahler deformations of
$(X,g)$ at fixed complex structure is
$H^{1,1}(X,\mathbb{R})$. Concretely, every $\delta J$ with
$\delta g_{i\bar j} = \mathrm{i}\,\delta J_{i\bar j}$ harmonic of type
$(1,1)$ deforms the metric to a nearby Ricci-flat one.
\end{proposition}

\begin{proposition}[Complex structure moduli; Kodaira--Spencer]\label{prop:CS-moduli}
The space of infinitesimal complex structure deformations of $X$ is
$H^1(X,T_X) \cong H^{2,1}(X)$, where the isomorphism is given by
contraction with the holomorphic volume form $\Omega$,
\begin{equation}
    \mu \;\mapsto\; \iota_{\mu}\Omega
        \;\in\; H^{2,1}(X) \,,
    \qquad
    (\iota_{\mu}\Omega)_{i j \bar k}
        \;=\; \Omega_{i j l}\, \mu^{l}{}_{\bar k} \,,
\end{equation}
where $\mu \in H^1(X,T_X)$ is the Beltrami differential, a $(0,1)$-form valued in $T^{1,0}X$.
\end{proposition}

We write $J=t^aD_a$, $a=1,\ldots,h^{1,1}$, for real K\"ahler
coordinates in a chosen divisor basis, with admissibility meaning
$J\in\KC(X)$, and $z^i$, $i=1,\ldots,h^{2,1}$, for complex
parameters on $H^{2,1}(X)$.
These are the precise geometric realisations of the schematic
parameters $u^i(x)$ in \S\ref{ssec:L1-framing}: they are
infinitesimal metric deformations $\delta g$ for which the deformed
metric remains Ricci-flat to first order. Yau's theorem upgrades
these infinitesimal deformations to genuine nearby Ricci-flat
K\"ahler metrics inside the K\"ahler cone and complex structure
moduli space.



\subsubsection{How ten-dimensional fields become four-dimensional fields}

We can now turn to the idea of Kaluza--Klein reductions.
Suppose first that $A=0$ and that the
internal space $X$ is fixed. A ten-dimensional scalar field admits
a formal expansion in eigenfunctions of the scalar Laplacian on
$X$,
\begin{equation}
    \Phi(x,y)=\sum_i \varphi_i(x) f_i(y),
    \qquad
    -\Delta_X f_i=\lambda_i f_i .
\end{equation}
After inserting this expansion into the ten-dimensional kinetic
term and integrating over $X$, the coefficient
$\varphi_i(x)$ becomes a four-dimensional scalar field with mass
squared proportional to $\lambda_i$, up to the conventional length
scale of $X$ and to corrections from warping or background fluxes.
Thus harmonic modes, $\lambda_i=0$, give massless fields before
additional interactions are included, while non-zero eigenvalues give
the Kaluza--Klein tower. A field is called \emph{light} if its mass
lies below the cutoff of the four-dimensional effective theory under
consideration.

The same idea applies to differential forms. If $C_p$ is a
$p$-form field and $\omega_\alpha$ is a harmonic $p$-form on
$X$, then
\begin{equation}
    C_p(x,y)=\sum_\alpha a_\alpha(x)\,\omega_\alpha(y)+\cdots,
    \qquad [\omega_\alpha]\in H^p(X,\mathbb{R}),
\end{equation}
produces scalar fields $a_\alpha(x)$ in four dimensions. This is
the first reason cohomology appears in the physics: it organises the
massless topological modes.

\subsection{Energy densities from moduli fields}

The previous subsection explained how cohomology classes and
families of internal metrics give rise to light four-dimensional
fields.  The next question is dynamical: once these fields have been
identified, what energy is assigned to a configuration of them?  This
is the point at which the geometry of the compact space becomes an
input to four-dimensional cosmology.

For the physics, the crucial extra step is to allow these parameters
to depend on the four-dimensional coordinates $x^\mu$.  A path in the
moduli space of Ricci-flat metrics then becomes a four-dimensional
field configuration.  If the fields vary in time, the internal metric
$g^{(6)}(y;\phi(t))$ also changes in time, and this time-dependence
contributes to the four-dimensional stress-energy.  In this sense, the
variation of the geometry of $X$ leads directly to cosmological
evolution in $3+1$ dimensions.

The phrase \emph{cosmological evolution} will be used here in a
minimal sense.  It means a solution of the four-dimensional Einstein
equations coupled to scalar fields, typically with a time-dependent
metric on $\mathcal{M}_4$ and time-dependent scalar fields
$\phi^A(t)$.  From the ten-dimensional point of view, this same
solution can be read as a family of internal geometries
$g^{(6)}(y;\phi(t))$ evolving over four-dimensional time.  Thus the
dynamical evolution of the internal metric is not an extra metaphor:
it is precisely the four-dimensional dynamics of the moduli fields.

This brings the discussion back to the central aim of the lecture
series.  Moving in the moduli space of internal geometries is, in the
four-dimensional effective theory, the same as changing the values of
scalar fields.  The scalar potential on this field space is the object
whose critical values are the vacuum energies that
Lectures~2 and~3 compute from toric data.


The metric itself also reduces. A family of Ricci-flat metrics
$g^{(6)}(y;\phi)$ depending on parameters $\phi^A$ gives
four-dimensional scalar fields by allowing those parameters to vary
over spacetime,
\begin{equation}
    g^{(6)}_{mn}(y;\phi)
    \quad\leadsto\quad
    g^{(6)}_{mn}\bigl(y;\phi(x)\bigr).
\end{equation}
If varying $\phi^A$ preserves the leading-order ten-dimensional
equations, then $\phi^A$ is a \emph{modulus}. 
For Calabi--Yau compactifications, the most important
examples are the K\"ahler moduli and complex structure moduli counted
by $h^{1,1}(X)$ and $h^{2,1}(X)$.  They are geometric fields in four
dimensions: their values determine the shape and size data of the
internal Calabi--Yau metric. 
Classically, these moduli have no potential:
moving in the moduli space still gives a vacuum solution, so
\begin{equation}
    V(\phi)=0
    \qquad
    \text{along the classical Calabi--Yau moduli space.}
\end{equation}


\begin{figure}[t!]
\begin{center}
\centering
\begin{tikzpicture}[scale=1.20,>=Latex]
  \begin{scope}
    \node at (1.75,2.95) {\revision{$T_{MN}=0$}};
    \draw[->,thin] (0,0) -- (0,2.25) node[above] {$V(\phi)$};
    \draw[->,thin] (0,0) -- (3.25,0) node[below right] {$\phi^A$};
    \draw[thick] (0.25,0.45) -- (2.95,0.45);
    \node[left] at (0,0.45) {$0$};
  \end{scope}

  \begin{scope}[xshift=5.45cm]
    \node at (1.75,2.95) {\revision{$T_{MN}\neq 0$}};
    \draw[->,thin] (0,0) -- (0,2.25) node[above] {$V(\phi)$};
    \draw[->,thin] (0,0) -- (3.25,0) node[below right] {$\phi^A$};
    \draw[thick]
      (0.25,0.65)
      .. controls (0.70,1.85) and (1.10,0.20) .. (1.55,0.95)
      .. controls (1.95,1.55) and (2.30,0.50) .. (2.75,1.35)
      .. controls (2.95,1.70) and (3.05,1.78) .. (3.15,1.85);
    \fill (1.48,0.92) circle (2pt);
  \end{scope}
\end{tikzpicture}
\end{center}
\vspace*{-0.5cm}
\caption{Classically the Calabi--Yau moduli are
flat directions: in the leading source-free approximation the
ten-dimensional stress-energy contribution is $T_{MN}=0$, and the
four-dimensional potential vanishes along the moduli space.  Fluxes,
localised sources, and further corrections are represented
schematically by $T_{MN}\neq0$ in the ten-dimensional equations; in
the four-dimensional theory they lift the flat moduli space to an
energy landscape $V(\phi)$.}
\end{figure}


Later we will include ingredients with non-trivial
stress-energy, such as fluxes, branes, orientifold planes, and
quantum corrections.  In ten dimensions these appear as
$T_{MN}\neq 0$ in \eqref{eq:EOMRvac}; in the four-dimensional
effective theory they generate a scalar potential $V(\phi)$ for the
fields $\phi^A$.
At the level of the resulting
four-dimensional effective action one has schematically
\begin{equation}
    S_4
    =
    \int_{\mathcal{M}_4}\mathrm{d}^4x\sqrt{-g_4}
    \left[
        \frac{\Pl^2}{2}R^{(4)}
        -G_{AB}(\phi)\,\partial_\mu\phi^A\partial^\mu\phi^B
        -V(\phi)
        +\cdots
    \right].
\end{equation}
For constant scalar fields, the scalar equations of motion reduce to
\begin{equation}
    \partial_A V(\phi_*)=0 .
\end{equation}
The four-dimensional Einstein equation then relates the value of the
potential to the spacetime curvature,
\begin{equation}
    R^{(4)}_{\mu\nu}
    =
    \frac{V(\phi_*)}{\Pl^2}\,g^{(4)}_{\mu\nu}
    \quad
    \text{in these conventions}.
\end{equation}
This is the energy-landscape picture used throughout the lectures:
geometry supplies fields and functions; a vacuum is a critical point;
and its vacuum energy is the critical value $V(\phi_*)$.
Lecture~2 identifies the geometric and arithmetic data
needed to write the four-dimensional functions $K$, $W$, and hence
$V$; Lecture~3 then uses these functions to compute $V(\phi_*)$ at
explicit candidate critical points.




\subsection{Why Calabi--Yau hypersurfaces in toric varieties?}\label{ssec:L1-toric}

The logic above tells us which manifolds
are admissible; it does not tell us which ones we can compute on.
The practical compactification programme is to enumerate large
families of Calabi--Yau threefold topologies, compute the
cohomological and intersection-theoretic data controlling the
four-dimensional fields, compute periods and instanton corrections
where needed, and then use these inputs to study potentials and
vacuum energies. For computability we use the toric construction of
Batyrev \cite{Batyrev:1993oya}, with which the audience is already
familiar. Projective hypersurfaces such as the quintic are familiar
first examples; Batyrev hypersurfaces in toric varieties are the
high-dimensional computational generalisation. A one-paragraph
summary, for the sake of fixing notation:

\begin{theorem}[Batyrev; CY threefolds from reflexive polytopes]\label{thm:batyrev}
Let $\Delta^\circ\subset N_\mathbb{R}$ be a reflexive $4$-polytope
(so $\Delta = (\Delta^\circ)^\vee \subset M_\mathbb{R}$ is also a
lattice polytope, and $\Delta^\circ$ has the origin as its unique
interior lattice point). Let $\Sigma_{\mathcal{T}}$ be the fan of a
fine, regular, star triangulation $\mathcal{T}$ of $\Delta^\circ$.
Then $V_{\Sigma_{\mathcal{T}}}$ is an MPCP toric ambient variety:
it is allowed to have terminal quotient singularities, but its
singular locus does not meet a generic anticanonical hypersurface in
the threefold case. Consequently the generic anticanonical
hypersurface
$X = \{f=0\}\subset V_{\Sigma_{\mathcal{T}}}$ is a smooth
Calabi--Yau threefold, and Batyrev gives explicit combinatorial
formulas for its Hodge numbers in terms of $\Delta,\Delta^\circ$
alone.
\end{theorem}

The Kreuzer--Skarke list \cite{Kreuzer:2000xy} of all $\Delta^\circ$
contains $473\,800\,776$ reflexive 4-polytopes; together with the
FRSTs $\mathcal{T}$ this gives a landscape of explicit Calabi--Yau
hypersurfaces in toric varieties covering
$1\leq h^{1,1}\leq 491$. This is the input to Lectures~2 and~3:
every concrete construction will be ``on a particular Batyrev
hypersurface in a specified toric ambient.''

Here is the toric punchline in the form needed for the rest of the
lectures. The Batyrev input $(\Delta^\circ,\mathcal{T})$ is not
the final object of interest. It is the starting point for a chain of
computations
\begin{equation}
    (\Delta^\circ,\mathcal{T})
    \leadsto X
    \leadsto
    \{\text{topological data, cones, periods, etc.}\}
    \leadsto
    (K,W,V)
    \leadsto
    V(\phi_*) \, .
\end{equation}
The last number is the vacuum energy of a four-dimensional critical
point. The list below is therefore a dictionary, not a collection of
unrelated toric facts. For each entry one should keep three questions
separate: what is the geometric object, where does it enter the
four-dimensional calculation, and how is it obtained from the toric
description?

\begin{itemize}[leftmargin=*]
    \item \emph{Wall data.}
	A torus-invariant divisor
    $\widehat D_\rho$ on the ambient toric variety restricts to a
    divisor $D_\rho=\widehat D_\rho\cap X$ on
    $X$, possibly empty for lattice points interior to facets.
    In favourable cases these restricted toric divisor classes span
    $H^{1,1}(X)$. Choosing a basis $D_a$, one computes
    \begin{equation}
        \kappa_{abc}=\int_X D_a\wedge D_b\wedge D_c,
        \qquad
        c_{2,a}=\int_X c_2(X)\wedge D_a .
    \end{equation}
    Together with the Hodge numbers $h^{1,1},h^{2,1}$, this completes the \emph{Wall data} of a Calabi--Yau threefold.
    For $J=t^aD_a$, these same numbers give the Calabi--Yau volume
    and divisor volumes
    \begin{equation}
        \begin{aligned}
        \Vol(X,J)
            &=
            \frac{1}{3!}\int_X J\wedge J\wedge J
            =
            \frac16\kappa_{abc}t^at^bt^c,\\
        \tau_a
            &=
            \frac12\int_{D_a} J\wedge J
            =
            \frac12\int_X D_a\wedge J\wedge J
            =
            \frac{\partial \Vol}{\partial t^a}
            =
            \frac12\kappa_{abc}t^bt^c .
        \end{aligned}
    \end{equation}

    \par\smallskip
    \noindent\emph{Used in:} the K\"ahler potential, the mirror prepotential, 
    and the non-perturbative superpotential.

    \par\noindent\emph{How:} In practice this data is obtained from the ambient
    toric intersection ring, the hypersurface class, and the
    Chern-class formula for toric hypersurfaces.  The implementation
    used in the computational parts of these lectures is
    \texttt{CYTools} \cite{Demirtas:2022hqf}; the same data are the
    Wall data used to compare toric phases in
    \cite{Gendler:2023ujl}.

    \item \emph{Cone data and flops.}
    The K\"ahler cone $\KC(X)$ records which real $(1,1)$-classes
    can be represented by K\"ahler forms; its dual Mori cone records
    effective curve classes. Curve volumes are the inequalities that
    keep the large-volume expansion under control. Different FRSTs
    can give adjacent K\"ahler chambers related by flops, and the
    union of these chambers is the extended K\"ahler cone
    $\eKC(X)$. In Lecture~3, moving through these adjacent chambers
    is the geometric part of the vacuum search: one follows the
    potential while checking that all relevant curve volumes remain
    positive and sufficiently large.

    \par\smallskip
    \noindent\emph{Used in:} the domain of validity of the large-volume
    expansion, the definition of K\"ahler chambers, and the geometric
    transitions used in the Lecture~3 search strategy.

    \par\noindent\emph{How:} \texttt{CYTools} computes Mori and K\"ahler cone data
    for toric phases and tracks the classical data across
    triangulations \cite{Demirtas:2022hqf}.  The secondary-fan and
    chamber-navigation layer used in the separate Lecture~3 and 4 discussion
    is provided by \texttt{fanroots} \cite{fanroots:code}.

    \item \emph{Divisor data.}
    Non-perturbative terms in the superpotential may come from ED3
    instantons or gaugino condensation on divisors. A necessary
    zero-mode test is expressed in terms of the Hodge numbers of the
    divisor, especially $h^{0,1}(D)$ and $h^{0,2}(D)$, or
    equivalently the arithmetic genus appearing in Witten's criterion.
    For prime toric divisors these Hodge numbers can often be read
    from the face of $\Delta$ dual to the lattice point defining
    the divisor. This is why a purely combinatorial computation can
    screen candidate instanton divisors before any four-dimensional
    potential is written.

    \par\smallskip
    \noindent\emph{Used in:} deciding which divisors can support D-instanton terms in $W_{\rm np}$.

    \par\noindent\emph{How:} For prime toric divisors,
    compute divisor Hodge numbers via the face-duality formulae, see e.g.
    \cite{Braun:2017nhi}.  This is a special toric shortcut.  For
    non-inherited or non-toric divisors, determining the relevant
    cohomology groups is a substantially harder sheaf-cohomology
    problem and should not be treated as automatic output of the
    toric combinatorics.

    \item \emph{Period data.}
    The flux superpotential in Type IIB is built from periods of the
    holomorphic three-form,
    \begin{equation}
        \Pi(z)=\left(\int_{\Gamma_I}\Omega(z)\right),
        \qquad \Gamma_I\in H_3(X,\mathbb{Z}) .
    \end{equation}
    Near a large complex structure point these periods are computed
    by mirror symmetry. In practice this means solving the
    Picard--Fuchs system of the mirror polytope and expressing the
    answer in flat coordinates. These periods determine
    $W_{\rm flux}$, and hence the flux value $W_0$ after the
    complex structure moduli and the axio-dilaton are stabilised.

    \par\smallskip
    \noindent\emph{Used in:} the flux superpotential $W_{\rm flux}$,
    in particular the stabilisation of complex structure moduli.

    \par\noindent\emph{How:} construct the GKZ /
    Picard--Fuchs system from the Batyrev mirror data, compute the
    Frobenius solutions near an LCS point, and then match the
    logarithmic solutions to an integral symplectic period basis using
    the large volume monodromy on the mirror side
    \cite{Hosono:1993qy,CoxKatz:1999,Demirtas:2023als}.

    \item \emph{Enumerative data.}
    Genus-zero Gopakumar--Vafa invariants $n_\beta$, for
    $\beta\in H_2(X,\mathbb{Z})$, are the integer invariants that
    repackage genus-zero Gromov--Witten invariants after the
    multiple-cover contributions have been separated
    \cite{Gopakumar:1998jq,Demirtas:2023als}. For the audience it is
    useful to remember two interpretations. Mathematically, they are
    the integral BPS/GV refinement of rational curve counts in class
    $\beta$. In the physical interpretation they are indexed BPS
    degeneracies, naturally described by M2-branes wrapping
    holomorphic curves in class $\beta$. They appear in the
    worldsheet-instanton part of the prepotential, so they correct
    the periods and can generate exponentially small terms used in
    small-$W_0$ constructions. Lecture~2 is devoted to making this
    computational mirror symmetry (CMS) / GV calculation explicit
    enough for the vacuum energy story.

    \par\smallskip
    \noindent\emph{Used in:} worldsheet-instanton corrections to the
    prepotential, convergence estimates for the instanton expansion,
    and the small terms that can enter small-$W_0$ and conifold-PFV
    constructions.

    \par\noindent\emph{How:} compare the instanton part of
    the mirror prepotential with the Frobenius period expansion, use
    the multiple-cover formula to pass from genus-zero
    Gromov--Witten coefficients to integral GV invariants, and organise
    the resulting finite computation by Mori-cone degree or by a
    decomposition-closed truncation
    \cite{Gopakumar:1998jq,Hosono:1993qy,Demirtas:2023als}.  In
    Lecture~3 the same curve-class data are then compared across
    flop-related chambers.
    
\end{itemize}


Many of the conceptual ingredients in this story are classical by now: Calabi--Yau compactification, mirror symmetry, periods, flux superpotentials, instanton corrections, and Gopakumar--Vafa invariants have been studied for decades. What has changed rather dramatically is our ability to compute these ingredients in practice, not only for a few hand-picked examples, but uniformly across large families with widely varying Hodge numbers. The main breakthrough for explicit string constructions is therefore algorithmic as much as conceptual: toric geometry, mirror symmetry, triangulation algorithms, period computations, and modern searches over flux and Kähler data turn the existence of a vast landscape into concrete, checkable Diophantine problems.


The companion \texttt{CYTools} session demonstrates the computational
pipelines on explicit examples: querying the Kreuzer--Skarke database,
constructing or sampling FRSTs, building a \texttt{CalabiYau} object,
computing intersection numbers and cone data, testing divisor
topology, and extracting truncated GV dictionaries. The remainder of
these lectures explains why these computations are the inputs to established string theory
constructions.


\paragraph{Lecture~1 punchlines.}
The first lecture should leave the following takeaways.
First, compactification turns internal geometry into
four-dimensional fields; the light fields are controlled by
cohomology, harmonic forms, and Ricci-flat metric deformations.
Second, uncorrected Calabi--Yau compactifications have flat moduli
spaces, so vacuum energies only become isolated numbers after
non-trivial sources of stress-energy generate a
potential. Third, Calabi--Yau threefolds enter because Yau's theorem
turns the nonlinear Ricci-flat metric problem into a computable
complex-geometric existence problem. Fourth, hypersurfaces a la Batyrev
make this computable at scale: toric data provide the intersection,
cone, divisor, period, and GV inputs that Lectures~2 and~3 feed into
the four-dimensional scalar potential.


That completes Lecture~1.

\iffalse
\begin{exercise}\label{ex:L1-1}
Let $X$ be a smooth Calabi--Yau threefold with $h^{1,1}=h^{2,1}=1$.
Compute $\chi(X)$ and write down the dimensions of all Hodge groups
$H^{p,q}(X)$.
\end{exercise}
\fi

\iffalse
\begin{exercise}\label{ex:L1-2}
Let $X$ be the resolution of a generic anti-canonical hypersurface
in a smooth $4$-dimensional toric variety $V_{\Sigma_{\mathcal{T}}}$ associated
to a reflexive polytope $\Delta^\circ$. Show that $\Omega \in
H^{3,0}(X)$ can be written explicitly as a residue of a meromorphic
form on the ambient space, and identify the K\"ahler cone of $X$ as
the (interior of the) image of the K\"ahler cone of $V_{\Sigma_{\mathcal{T}}}$
under restriction.
\end{exercise}
\fi
